\documentclass[10pt,a4paper]{amsart}
\usepackage[margin=19mm]{geometry}
\usepackage{amsmath,amssymb,mathtools}
\usepackage[scr=rsfs,cal=euler]{mathalfa}
\usepackage{xfrac}
\usepackage[unicode,hidelinks,bookmarks=false]{hyperref}
\newcommand{\ie}{i.e.\ }
\newcommand{\cf}{cf.\ }
\newcommand{\resp}{respectively\ }
\newcommand{\Subsection}{\subsection}
\providecommand{\nobreakdash}{\nobreak\hbox{-}}
\setlength{\emergencystretch}{2em}
\multlinegap0pt
\usepackage{bm}
\usepackage{bbm}
\usepackage{dsfont}
\usepackage{xspace}
\usepackage{cancel}
\usepackage{mathtools}
\usepackage{cleveref}
\usepackage[shortlabels]{enumitem}
\usepackage{amsthm}
\makeatletter
\@ifundefined{theorem}{\newtheorem{theorem}{Theorem}}{}
\@ifundefined{corollary}{\newtheorem{corollary}[theorem]{Corollary}}{}
\@ifundefined{lemma}{\newtheorem{lemma}[theorem]{Lemma}}{}
\@ifundefined{remark}{\newtheorem{remark}[theorem]{Remark}}{}
\makeatother
\newcommand{\B}{\mathbb B}
\newcommand{\R}{\mathbb R}
\newcommand{\sB}{{\mathcal B}}
\title[Well-posedness of Lur'e systems with feedthrough]{Well-posedness of Lur'e systems with feedthrough}
\author{Chris Guiver, Hartmut Logemann}
\date{Source: arXiv:2508.16221v1 (CC BY 4.0)}
\begin{document}
\maketitle

\noindent\emph{Source and licence.} Well-posedness of Lur'e systems with feedthrough. Version arXiv:2508.16221v1 (CC BY 4.0), distributed under the Creative Commons Attribution 4.0 International licence, as recorded in the arXiv licence metadata for this exact version and shown on the abstract page https://arxiv.org/abs/2508.16221v1. Full rights notice: licenses/arxiv-2508.16221-CC-BY-4.0.txt.\par\smallskip

\noindent\textit{Editorial scope.} Counted unit: the Corollary whose source label is \texttt{cor:inclusion\_wp} in the article (source0.tex lines 1577--1581, source label \texttt{cor:inclusion\_wp}), delivered together with the article's own complete original proof of it (source0.tex lines 1583--1621, one \texttt{proof} environment of 39 lines). No mathematics is written, repaired or re-derived: statement and proof are the article's own text line for line. Required context retained uncounted: eq:lure\_ss (lines 231--237); thm:single\_valued (lines 813--833); eq:single\_valued (lines 818--820); rem:a.e. (lines 838--857); eq:al\_bound (lines 885--887); ls:a1 (lines 1226--1234); ls:a3 (lines 1226--1234); thm:inclusion\_wp (lines 1238--1252); ls:inclusion\_wp\_s\_2 (lines 1242--1251); rem:a.e.' (lines 1258--1275); ls:a1' (lines 1261--1268); ls:a2' (lines 1261--1268); ls:a3' (lines 1261--1268); lem\_a:gronwall (lines 1789--1796). Every definition, display and label of those lines is retained unchanged. Omitted from this package and not redistributed: 1--230, 238--812, 821--837, 858--884, 888--1225, 1235--1237, 1252--1257, 1269--1576, 1582--1582, 1622--1788, 1797--2060. Nothing inside a retained excerpt is omitted or altered: every retained body is delivered exactly as the article's lines are written, so no omission receipt of any kind accompanies this package. Citations: the retained excerpts carry no citation command at all, so no bibliography entry of the article is transcribed and no reference is redistributed. Reference adaptation: none was needed and none was made; every reference inside the delivered unit resolves inside it, so no unresolved reference or citation is produced. Printed slips kept literal: no printed word, symbol, hypothesis or number of the counted statement or of its proof was altered. Layout and class: the article's own class is \texttt{article} with packages including xcolor, amsmath, amsthm, amssymb, bbm, enumitem, epsfig, graphics, graphicx, fancyhdr, geometry, parskip, cleveref, tikz; the wrapper is the lane's pinned \texttt{amsart} editorial wrapper at 10pt on A4, which restates the theorem-like environments the retained text needs and adds the article's own macro definitions that the retained text needs (\texttt{\textbackslash B}, \texttt{\textbackslash R}, \texttt{\textbackslash sB}), with extra wrapper packages (bm, bbm, dsfont, xspace, cancel, mathtools, cleveref, enumitem). The delivered numbering is the unit's own. Assets: no retained span contains a figure environment, a graphics-inclusion command, a PDF-page inclusion, a table or a TikZ picture; no image file is copied into this package, the package declares no asset dependency and no image file is redistributed with it. Licence: version arXiv:2508.16221v1 of this article is distributed under the Creative Commons Attribution 4.0 International licence, as recorded in the arXiv licence metadata for this exact version and shown on the abstract page https://arxiv.org/abs/2508.16221v1. A full rights notice is delivered at licenses/arxiv-2508.16221-CC-BY-4.0.txt. Characters: every retained excerpt is pure ASCII, so the delivered engine, XeLaTeX with its default Latin Modern text font, reports no missing character.
 \begin{subequations}\label{eq:lure_ss}
 \begin{align}
 \dot x(t) & = Ax(t)+ Bf(t,y(t)) + B_{\rm e}v(t),\quad x(t_0)=x^0,\,\,\,  t\geq t_0\geq 0,
 \label{eq:lure_ss_a}\\
 y(t) & = Cx(t)+Df(t,y(t))+ D_{\rm e}v(t),\label{eq:lure_ss_b}
 \end{align}
 \end{subequations}

\begin{theorem}\label{thm:single_valued}
Assume that $F_t$ is locally injective for every $t\geq 0$ and that
$F_t$ is radially unbounded, locally uniformly in $t$. Furthermore,
assume that, for every compact subset $K\subset\R^p$, there exists
$\varphi\in L^1_{\rm loc}(\R_+,\R_+)$ such that
\begin{equation}\label{eq:single_valued}
\sup_{\xi\in K}\|f(t,\xi)\|\leq\varphi(t)\quad\forall\,t\geq 0.
\end{equation}
Let $v\in L^\infty_{\rm loc}(\R_+,\R^{m_{\rm e}})$, $t_0\geq 0$ and $x^0\in\R^n$. The following statements hold.
%
\begin{enumerate}[label = {\bf (\arabic*)}, ref={\rm (\arabic*)}, leftmargin = 0ex, itemindent = 4.6ex]
\item\label{ls:single_valued_s_1}
There exists $(v,x,y)\in\tilde\sB(S^f,t_0)$ such that $x(t_0)=x^0$.
%
\item\label{ls:single_valued_s_2}
If $(v,x,y)\in\tilde\sB(S^f,t_0)$ with bounded maximal interval of definition $[t_0,\tau)$,
where $t_0<\tau<\infty$, then $\int_{t_0}^\tau\|Bf(t,y(t))\|{\rm d}t=\infty$, $y\not\in L^\infty([0,\tau),\R^p)$ and $\|x(t)\|\to\infty$ as $t\uparrow\tau$. In particular,
system \eqref{eq:lure_ss} has the blow-up property.
%
\end{enumerate}
\end{theorem}

\begin{remark}\label{rem:a.e.}
\begin{rm}
The above result remains valid under slightly weaker assumptions on $F$. Assume that \eqref{eq:single_valued} holds, $F_t$ is locally injective for almost every $t\geq 0$
and $F_t$ is radially unbounded, {\it locally essentially uniformly in $t$}, in the sense
that, for all $\rho>0$ and all compact $T\subset\R_+$, there exists $\sigma\geq 0$ such that $\|F_t(\xi)\|\geq\rho$ for almost every $t\in T$ and all $\xi$ with $\|\xi\|\geq\sigma$.
We claim that the conclusions of \Cref{thm:single_valued}
are still valid. To see this, note that there exists a null set $E\subset\R_+$ such that
$F_t$ is locally injective for every $t\in\R_+\backslash E$, and, for all $\rho>0$ and all compact $T\subset\R_+$, there exist $\sigma\geq 0$ such that $\|F_t(\xi)\|\geq\rho$ for all $t\in T\backslash E$ and all $\xi$ with $\|\xi\|\geq\sigma$. Defining
\[
\hat f(t,\xi):=\left\{ \begin{aligned} & f(t,\xi), & & (t,\xi)\in(T\backslash E)\times\R^p, \\
& 0, & & (t,\xi)\in E\times\R^p
\end{aligned}\right.\qquad\mbox{and}\qquad
\hat F_t(\xi):=\hat F(t,\xi):=\xi-D\hat f(t,\xi),
\]
it is clear that $\hat f$ is a Caratheodory function, $\tilde\sB(S^{\hat f},t_0)=
\tilde\sB(S^f,t_0)$, $\hat F_t$ is locally injective
for every $t\geq 0$ and $\hat F_t$ is radially unbounded, locally uniformly in $t$. Hence, the above claim follows from an application of \Cref{thm:single_valued} to system \eqref{eq:lure_ss} with $f$ replaced by $\hat f$.
\hfill$\Diamond$
\end{rm}
\end{remark}

\begin{equation}\label{eq:al_bound}
\sup_{\|\xi\|\geq\rho}\frac{\|f(t,\xi)\|}{\|\xi\|}\leq c\qquad\mbox{for a.e. $t\in T$}.
\end{equation}

\begin{enumerate}[label = {\bfseries (A\arabic*)}, ref = {\rm (A\arabic*)}]
%
\item \label{ls:a1} For every $t\geq 0$, the set~$F_t^{-1}(\xi)=\{\zeta\in\R^p:F_t(\zeta)=\xi\}$ is nonempty for all $\xi\in{\rm im}\,(C,D_{\rm e})$.
%
\item \label{ls:a2} $F_t$ is radially unbounded, locally uniformly in $t$.
%
\item \label{ls:a3} For every~$t\geq 0$ and every~$\xi\in {\rm im}\,(C,D_{\rm e})$, the set~$f(t,F_t^{-1}(\xi))=\{f(t,\zeta):\zeta\in F_t^{-1}(\xi)\}$ is convex.

\end{enumerate}

\begin{theorem}\label{thm:inclusion_wp}
Assume
that \ref{ls:a1}--\ref{ls:a3} are satisfied and, for all compact sets $K\subset\R^p$, there exists $\varphi\in L^1_{\rm loc}(\R_+,\R_+)$ such that \eqref{eq:single_valued} holds.
Let $v\in L^\infty_{\rm loc}(\R_+,\R^{m_{\rm e}})$, $t_0\geq 0$ and $x^0\in\R^n$. The following statements hold.
\begin{enumerate}[label = {\bf (\arabic*)}, ref={\rm (\arabic*)}, leftmargin = 0ex, itemindent = 4.6ex]
%
\item \label{ls:inclusion_wp_s_1} There exists $(v,x,y)\in\tilde\sB(S^f,t_0)$ such that $x(t_0)=x^0$.
%
\item \label{ls:inclusion_wp_s_2}
If $(v,x,y)\in\tilde\sB(S^f,t_0)$ has a bounded maximal interval of definition $[t_0,\tau)$,
where $t_0<\tau<\infty$, then $\int_{t_0}^\tau\|Bf(t,y(t))\|{\rm d}t=\infty$, $y\not\in L^\infty([0,\tau),\R^p)$ and $\|x(t)\|\to\infty$ as $t\uparrow\tau$. In particular,
system \eqref{eq:lure_ss} has the blow-up property.
%
\end{enumerate}
\end{theorem}

\begin{remark}\label{rem:a.e.'}
\begin{rm}
Consider the following assumptions which are slightly weaker than \ref{ls:a1}-\ref{ls:a3}.
\begin{enumerate}[label = {\bfseries (B\arabic*)}, ref = {\rm (B\arabic*)}]
%
\item \label{ls:a1'} For almost every $t\geq 0$, the set~$F_t^{-1}(\xi)=\{\zeta\in\R^p:F_t(\zeta)=\xi\}$ is nonempty for all $\xi\in{\rm im}\,(C,D_{\rm e})$.
%
\item \label{ls:a2'} $F_t$ is radially unbounded, locally essentially uniformly in $t$, in the sense of \Cref{rem:a.e.}.
%
\item \label{ls:a3'} For almost every~$t\geq 0$ and every~$\xi\in {\rm im}\,(C,D_{\rm e})$, the set~$f(t,F_t^{-1}(\xi))=\{f(t,\zeta):\zeta\in F_t^{-1}(\xi)\}$ is convex.
\end{enumerate}
We claim that \Cref{thm:inclusion_wp} remains valid if assumptions \ref{ls:a1}-\ref{ls:a3} are replaced by \ref{ls:a1'}-\ref{ls:a3'}. Indeed, if \ref{ls:a1'}-\ref{ls:a3'} hold, then
there exists a null set $E\subset\R_+$ such that the  Caratheodory functions $\hat f$ and $\hat F$ defined in \Cref{rem:a.e.} satisfy \ref{ls:a1}-\ref{ls:a3}. Consequently, \Cref{thm:inclusion_wp} applies to system \eqref{eq:lure_ss} with $f$ replaced by $\hat f$
and the claim follows from the identity  $\tilde\sB(S^{\hat f},t_0)=
\tilde\sB(S^f,t_0)$.
\hfill$\Diamond$
\end{rm}
\end{remark}

\begin{lemma}\label{lem_a:gronwall}
Let $0\leq t_0<\tau\leq\infty$, $c\geq 0$, $h\in L^p_{\rm loc}(\R_+,\R)$ be non-negative and
$y\in L^q_{\rm loc}(\R_+,\R)$ , where $1\leq p,q\leq\infty$ are such that $1/p+1/q=1$. If
\[
y(t)\leq c+\int_{t_0}^t h(s)y(s){\rm d}s,\quad\mbox{for a.e. $t\in[t_0,\tau)$},
\]
then $y(t)\leq ce^{\int_{t_0}^t h(s){\rm d}s}$ for almost every $t\in[t_0,\tau)$.
\end{lemma}

\begin{corollary}\label{cor:inclusion_wp}
Assume that \ref{ls:a1} and \ref{ls:a3} hold, for all compact sets $K\subset\R^p$, there exists $\varphi\in L^1_{\rm loc}(\R_+,\R_+)$ such that \eqref{eq:single_valued} holds, and,
for all compact sets $T\subset\R_+$, there exist $\rho>0$ and $c>0$ such that $c\|D\|<1$ and
\eqref{eq:al_bound} is satisfied. Then $\tilde\sB(S^f,t_0)=\sB(S^f,t_0)$ for all $t_0\geq 0$ {\rm (}that is, \eqref{eq:lure_ss} is forward complete{\rm )}. Furthermore, for all $v\in L^\infty_{\rm loc}(\R_+,\R^{m_{\rm e}})$, $t_0\geq 0$ and $x^0\in\R^n$, there exists $(v,x,y)\in \sB(S^f,t_0)$ such that $x(t_0)=x^0$.
\end{corollary}

\begin{proof}
As $c\|D\|<1$, it follows from \eqref{eq:al_bound} that $F_t$ is radially unbounded, locally essentially uniformly, in the sense of \Cref{rem:a.e.}, that is, \ref{ls:a2'} in \Cref{rem:a.e.'}
is satisfied. Hence, by  \Cref{rem:a.e.'}, the conclusions of \Cref{thm:inclusion_wp} hold. Therefore, we only need to establish that  $\tilde\sB(S^f,t_0)=\sB(S^f,t_0)$ for all $t_0\geq 0$. To this end, let $(v,x,y)$ be a pre-trajectory defined on $[t_0,\tau)$, where $t_0<\tau<\infty$. It is sufficient to prove that $y\in L^\infty([t_0,\tau),\R^p)$. Indeed, in this case,
statement \ref{ls:inclusion_wp_s_2} of \Cref{thm:inclusion_wp} implies that every maximally defined trajectory starting at $t_0$ is defined on $[t_0,\infty)$ and is therefore in $\sB(S^f,t_0)$.

To show that $y\in L^\infty([t_0,\tau),\R^p)$, note that, by \eqref{eq:al_bound}, there exist $\rho>0$, $c>0$ and a null set $E\subset[t_0,\tau]$
such that $c\|D\|<1$ and
\begin{equation}\label{eq:al_bound'}
\|f(t,\xi)\|\leq c\|\xi\|\quad\forall\,t\in [t_0,\tau]\backslash E,\,\,\forall\,\xi\in\R^p\backslash\B(0,\rho).
\end{equation}
For $t\in [t_0,\tau]\backslash E$, if $\|y(t)\|\geq\rho$, we have
\[
\|y(t)\|\leq\frac{1}{1-c\|D\|}\|Cx(t)+D_{\rm e}v(t)\|,
\]
and thus,
\[
\|y(t)\|\leq\rho+\frac{1}{1-c\|D\|}\|Cx(t)+D_{\rm e}v(t)\|\quad\mbox{for a.e. $t\in[t_0,\tau)$.}
\]
An application of the variation-of-parameters formula for $x$ yields that there
exist constants $c_1,c_2>0$ and a continuous non-negative function $h$ defined
on $[t_0,\tau]$ such that
\begin{equation}\label{eq:y_inequality}
\|y(t)\|\leq c_1+c_2\int_{t_0}^th(s)\|f(s,y(s))\|{\rm d}s\quad\mbox{for a.e. $t\in[t_0,\tau)$.}
\end{equation}
By \eqref{eq:single_valued}, there exists $\varphi\in L^1_{\rm loc}(\R_+,\R_+)$ such that
\[
\sup_{\xi\in\overline{\B(0,\rho)}}\|f(t,\xi)\|\leq\varphi(t)\quad\mbox{for a.e. $t\geq 0$}.
\]
Combining this with \eqref{eq:al_bound'} shows that
\[
\|f(t,y(t))\|\leq\varphi(t)+c\|y(t)\|\quad\mbox{for a.e. $t\in[t_0,\tau)$},
\]
which in conjunction with \eqref{eq:y_inequality} leads to
\[
\|y(t)\|\leq c_3+c_4\int_{t_0}^th(s)\|y(s)\|{\rm d}s\quad\mbox{for a.e. $t\in[t_0,\tau)$,}
\]
where $c_3$ and $c_4$ are suitable positive constants. An application of the Gronwall lemma (see \Cref{lem_a:gronwall} in the Appendix for a suitably general version) yields that $\|y(t)\|\leq c_3+\exp\big(c_4\int_{t_0}^t h(s){\rm d}s\big)$ for almost every $t\in [t_0,\tau)$, whence
$y\in L^\infty([t_0,\tau),\R^p)$.
\end{proof}

\end{document}
